\chapter{Metodología de selección y fuentes}
\label{chap:metodologia}
\section{Criterios de inclusión y exclusión}
\pendiente{Definir relevancia, actualidad, profundidad técnica, impacto social y disponibilidad de evidencia. Documentar candidatos descartados y evitar contar varias noticias del mismo hecho como casos diferentes.}
\section{Estrategia de búsqueda}
\pendiente{Registrar consultas, repositorios, idiomas, fechas de búsqueda y criterios de cierre. Priorizar avisos del fabricante, CERT/CSIRT, registros de vulnerabilidades, informes oficiales y trabajos originales.}
\section{Evaluación y contraste de fuentes}
\pendiente{Identificar autoría, fecha, versión, proximidad al hecho y posibles intereses. Contrastar afirmaciones relevantes con evidencia independiente; varias noticias que reproducen el mismo comunicado no son fuentes independientes.}
\section{Tratamiento técnico y CVSS}
\pendiente{Distinguir incidente, vulnerabilidad y amenaza. Registrar CVE cuando exista. Para CVSS, anotar versión, vector, puntuación, emisor y fecha; distinguir valores publicados de estimaciones propias. Si falta un score y el caso es evaluable, intentar calcularlo justificando cada métrica y sus supuestos. Si no aplica, explicar por qué; no inventar CVE ni forzar una puntuación.}
\section{Trazabilidad, limitaciones e integridad académica}
\pendiente{Vincular cada afirmación verificable con su fuente. Separar hechos, estimaciones e interpretación propia. Registrar contradicciones, incertidumbres y límites de acceso. Declarar herramientas de apoyo y su uso conforme a las normas de la asignatura.}
